% Chapter 2 -- Programming in Turbo Initech (model: TP 6.0 UG Ch 2, AN-19:
% continuous tutorial prose, milestone A-heads, short complete programs
% typed, run and dissected, no exercises; ends on a forward pointer).
% Every program and its output is replayed by test-pubs-turbo (Free Pascal,
% the seed, Turbo Initech bare image, and InitechDOS).  Language rules as
% defined in the Programmer's Guide (ADR-0007 DEC-02 as tps.pas implements
% it).
\Gates{test-compiler, test-compiler-os, test-pubs-turbo}
\chapter{Programming in Turbo Initech}
\ix{tutorial}

This chapter teaches the Turbo Initech language through a series of short
programs.  Each one adds one idea.  Type each program, compile it as
Chapter 1 shows, and run it; the output printed after each program is
exactly what it writes on COM1.

\section{Numbers and arithmetic}
\ix{arithmetic}
The first real program computes a week's pay:

\Program{pay}

\ProgramOutput{pay}

\SideHead{Comments}The text between braces is a comment; the compiler
ignores it.

\SideHead{Constants}The \rw{const} section gives the name \id{RateCents} to
the value 1250.  A constant cannot change while the program runs.

\SideHead{Variables}The \rw{var} section declares two \term{variables} of
type \rw{integer}.  Every variable must be declared before it is used.

\SideHead{Assignment}\code{Gross := Hours * RateCents} computes the product
and stores it in \id{Gross}.  The operator \code{:=} is read
``becomes.''

\SideHead{Division}Turbo Initech has whole numbers only.  \rw{div} divides
and discards the remainder, and \rw{mod} gives the remainder, so
\code{Gross div 100} is the dollars and \code{Gross mod 100} the cents.
The last line prints the two digits of the cents one at a time, so that 5
cents prints as \code{05}.

\MarginNote{An integer holds about two billion.}Integers in Turbo Initech
range from $-2147483648$ to 2147483647, so amounts of money up to about
21 million dollars can be kept in cents.

\section{Making decisions}
\ix{if statement}
The \rw{if} statement chooses between two actions:

\Program{overtime}

\ProgramOutput{overtime}

The condition \code{Hours > Normal} is either \rw{true} or \rw{false}.  When
it is true the statements between \rw{begin} and \rw{end} after \rw{then}
are carried out; otherwise the statement after \rw{else}.  A
\rw{begin}\dots\rw{end} pair groups several statements into one.  Note
that there is no semicolon before \rw{else}.

\section{Repeating actions}
\ix{loops}
\ix{for statement}
\ix{while statement}
\Program{table}

\ProgramOutput{table}

The \rw{for} statement counts: the outer loop runs \id{Row} from 1 to 4, and
for each row the inner loop runs \id{Col} from 1 to 4.
\std{write}{Write} writes without ending the line, and the
\std{writeln}{Writeln} with no arguments ends it.  The \rw{while} statement
repeats as long as its condition is true; here it counts down.  A third
loop, \rw{repeat} \dots{} \rw{until}, tests its condition at the end, so its
statements always run at least once.

\section{Characters and strings}
\ix{strings}
\ix{characters}
\Program{initials}

\ProgramOutput{initials}

A \rw{string} holds up to 255 characters; \code{Name[I]} is its
\id{I}th character, a value of type \rw{char}.  \std{length}{Length} gives
the number of characters, and \code{+} joins strings and characters.  The
Boolean variable \id{NewWord} remembers whether the next letter begins a
word.

\section{Arrays}
\ix{arrays}
An \term{array} holds a numbered series of values of one type:

\Program{sales}

\ProgramOutput{sales}

\code{Units[Q]} is the element number \id{Q}.  The bounds of the array,
1 to \id{Quarters}, are part of its declaration.  Turbo Initech does not
check that an index is within the bounds when the program runs, so keep
your loops within them.

\section{Records}
\ix{records}
A \term{record} groups values of different types under one name:

\Program{staff}

\ProgramOutput{staff}

The type \id{Person} is declared in a \rw{type} section.  \code{Team[I].Desk}
is the field \id{Desk} of the \id{I}th person.  Fields may be integers,
Boolean values or characters.

\section{Procedures and functions}
\ix{procedures}
\ix{functions}
A program grows clearer when it is divided into routines:

\Program{tax}

\ProgramOutput{tax}

A \term{function} computes a value: \id{Withholding} returns one quarter or
one fifth of its argument, by assigning to its own name.  A
\term{procedure} carries out an action: \id{Report} prints an amount.
\id{Tag} and \id{Cents} are \term{parameters}; each call supplies values for
them.  A call always has parentheses, even when there are no arguments.

\section{Formatting output}
\ix{columns}
\ix{field width}
\std{write}{Write} has no field widths: it writes each value in as many
characters as it needs.  To line numbers up in columns, write the spaces
yourself.  \id{WriteRight} counts the digits of a number and writes enough
spaces in front of it to fill a field:

\Program{columns}

\ProgramOutput{columns}

\section{Functions that answer questions}
\ix{Boolean functions}
A function of type \rw{boolean} turns a test into a name.  \id{IsLeap}
applies the calendar rules in order, so that a later rule corrects an
earlier one:

\Program{calendar}

\ProgramOutput{calendar}

\id{DaysIn} starts from the usual answer and changes it where the month
requires; one function may call another, provided the called function is
declared first.

\section{Recursion}
\ix{recursion}
A routine may call itself.  Moving a stack of disks from one peg to
another, one disk at a time and never a larger disk on a smaller one, is
three steps: move all but the largest disk out of the way, move the largest
disk, and move the others back on top of it.  The first and last steps are
the same problem with one disk fewer:

\Program{hanoi}

\ProgramOutput{hanoi}

Each call has its own parameters, so the four parameters of the outer call
are unchanged when the inner calls return.

\section{Files}
\ix{files}
Programs keep data from one run to the next in files:

\Program{savefile}

\ProgramOutput{savefile}

\std{assign}{Assign} gives the file variable a file name;
\std{rewrite}{Rewrite} creates the file; \std{blockwrite}{BlockWrite}
writes the characters of the string; \std{reset}{Reset} reopens the file at
its beginning; and \std{blockread}{BlockRead} reads the characters back into
a string.  Because \std{blockread}{BlockRead} does not change the length of
the string, the program first makes \id{Back} as long as the text it will
read.

\section{Where to go from here}
You now know every kind of statement and declaration in Turbo Initech.
Chapter 3 puts them together in a larger program.
